% Generated draft: requires mathematical and source review.
\documentclass[11pt]{article}
\usepackage[margin=25mm]{geometry}
\usepackage{fontspec}
\setmainfont{Latin Modern Roman}
\usepackage{amsmath,amssymb,mathtools,unicode-math}
\setmathfont{Latin Modern Math}
\usepackage{enumitem,needspace,xurl,hyperref}
\setlength{\parindent}{0pt}
\setlength{\parskip}{5pt}
\setlength{\emergencystretch}{3em}
\newcommand{\N}{\mathbb{N}}
\newcommand{\Z}{\mathbb{Z}}
\newcommand{\Q}{\mathbb{Q}}
\newcommand{\R}{\mathbb{R}}
\newcommand{\C}{\mathbb{C}}
\newcommand{\F}{\mathbb{F}}
\newcommand{\A}{\mathbb{A}}
\newcommand{\PP}{\mathbb{P}}
\newcommand{\E}{\mathbb{E}}
\newcommand{\eps}{\varepsilon}
\newcommand{\dd}{\,\mathrm{d}}
\begin{document}
\section*{UnsolvedMath ID: 11000176}
\textbf{Source problem code:} catalogue entry
\subsection{Definitions and mathematical statement}
Let \(M_g\) denote the mapping class group in the source and let \(X\) be one of the complexes discussed in its unprovided \(\S1\).

\emph{Problem 3.5.} Is there a natural quotient complex of such an \(X\) useful for constructing nonfaithful representations of \(M_g\)? A possible precise setup is an \(M_g\)-equivariant quotient \(q:X\to Y\) inducing
\[
 \rho_Y:M_g\longrightarrow\operatorname{Aut}(Y)
\]
with nontrivial kernel, followed by any further construction required for the intended representation target. The source does not say whether ``representation'' means this combinatorial action, a linear representation, or another kind. It also gives no naturalness or usefulness criterion. Collapsing \(X\) to a point formally produces a trivial action, but does not by itself meet the stated research purpose; no such construction is asserted as an answer.

The accompanying discussion proposes using answers to Problems 3.1, 3.2, or 3.5 to obtain invariants distinguishing Heegaard splittings. It recalls a successful symplectic-representation invariant in [4] separating inequivalent minimal splittings, says later student attempts with other representations were unsuccessful, and proposes revisiting that direction. These are source history, not a theorem that other representations cannot work.

For a fixed genus and handlebody marking, let \(\mathcal H\leq M_g\) be the handlebody subgroup. The source identifies the equivalence class of a splitting \(H\cup_\varphi H'\) with the double coset
\[
 \mathcal H\varphi\mathcal H\in\mathcal H\backslash M_g/\mathcal H.
\]
A representation \(\rho\) correspondingly gives image double-coset data \(\rho(\mathcal H)\rho(\varphi)\rho(\mathcal H)\), but usefulness requires genuinely distinguishing intended splitting classes. The exact orientation, side-ordering, and handlebody-identification conventions are omitted. The source recalls that different splittings of one \(3\)-manifold can become equivalent after stabilization and asks broadly about the collection of all genera; it does not require the proposed splitting invariant to identify or distinguish classes after stabilization.
\subsection{Short English statement}
Seek a natural quotient of a surface complex that yields useful nonfaithful mapping-class representations. The source hopes to use such representations to distinguish Heegaard double cosets, following an earlier symplectic invariant, but omits the complexes and precise representation target.
\subsection{Sources}
\begin{itemize}
\item Local UnsolvedMath registry, permanent ID 11000176, source code catalogue entry, Problem 3.5, PDF page 161. No individual proposer is supplied; page distinguishes repeated local labels.
\item Farb, \emph{Problems on Mapping Class Groups and Related Topics}. Supplied URL: \url{https://www.math.uchicago.edu/~farb/papers/mcgbook.pdf}; recorded access 2026-07-29; extraction pdf-font-structure. The volume contains 329 explicitly typeset headings across contributed chapters.
\item The statement refers to its \(\S1\) complexes, Problems 3.1 and 3.2, and reference [4] for the earlier symplectic splitting invariant, none supplied in full.
\item Supplied status open and OPEN-TRIAGE. The research summary says only ``Cannot verify.'' Current primary-literature status and redistribution terms for the publicly accessible author/publisher copy remain NEEDS\_REVIEW.
\end{itemize}
\end{document}
